\subsection{A countable target for presentations with two new variables per relator}

\begin{claim}{A.25}{Proved}{Sparse relative presentations admit small coefficient-preserving images}
Let $\nu$ be an infinite cardinal and let $H$ be a group with
$|H|\leq\nu$. There are a group $E_H$ with $|E_H|\leq\nu$ and an
embedding $j:H\hookrightarrow E_H$ with the following property.
Consider any relative presentation
\[
 G=(H*F(S))/\langle\!\langle w_r:r\in R\rangle\!\rangle,
 \qquad w_r\in H*F(S),
 \tag{A.25.1}
\]
whose exponent-sum homomorphism
\[
 \delta:\mathbb Z^{(R)}\longrightarrow\mathbb Z^{(S)}
 \tag{A.25.2}
\]
is an isomorphism. Suppose each word $w_r$ involves at most two
distinct variables from $S$, and the exponent sum of each variable
that occurs in it is nonzero. Then there is a homomorphism
$\Phi:G\to E_H$ whose composite with the canonical map $H\to G$
is exactly $j$. In particular this canonical map is injective.
The group $E_H$ is fixed for $H$, independently of $S$, $R$, and
their cardinalities.

Here the support condition concerns the full chosen word, including
all variable occurrences; it is not a condition only on the nonzero
entries of (A.25.2).

Consequently, let $X=BG$ be an actual two-dimensional classifying
CW complex, with one vertex and finite edge-word attaching maps.
Assume $H_1(G;\mathbb Z)$ and $H_2(G;\mathbb Z)$ are countable.
Suppose there are countable sets $S_*$ of one-cells and $R_*$ of
two-cells such that every attaching word outside $R_*$ has at
most two distinct variables outside $S_*$, each of nonzero
exponent sum. Then every nonidentity element of $G$ survives
in a countable quotient of $G$.
\end{claim}

\begin{proof}
Fix $j:H\hookrightarrow E_H$ by the finite-equation construction
in A.19. Explicitly, perform countably many passes; during each
pass enumerate all at most $\nu$ finite systems over the previous
group with rationally independent exponent rows, and successively
adjoin their universal solutions. The injectivity theorem
\cite[Theorem~1]{D_OpenAI2026Nonsingular} makes all coefficient
maps injective. At limit stages take group colimits. The result
has size at most $\nu$, contains $H$, and solves every finite
nonsingular system over itself. We will solve all the equations
in (A.25.1) inside this group, without taking a reduced power.

\paragraph{The graph forced by the specified exponent map.}
Make a multigraph with vertex set $S$. A relator involving two
distinct variables gives an edge between them. A relator
involving one variable gives a ground mark at that vertex.
Multiple edges and multiple ground marks are retained. There
are no relators involving zero variables, because their zero
columns would contradict injectivity of $\delta$.

Every connected component of this multigraph has at most one
cycle or one ground mark in total. To prove this, suppose a
component contains two cycles, a cycle and a ground mark, or
two ground marks. These data and paths connecting them lie
in a finite connected subgraph. Choose such a subgraph with
$v$ vertices, $e$ ordinary edges, and $b$ ground marks, so
that its cycle rank plus $b$ is at least two. Since its cycle
rank is $e-v+1$, it has
\[
 e+b\geq v+1.
\]
The corresponding $e+b$ columns of $\delta$ are supported
on those $v$ vertices. They are linearly dependent over
$\mathbb Q$, and clearing denominators contradicts
injectivity of $\delta$. This also treats parallel edges:
two parallel edges constitute a cycle of length two.

A component cannot be a tree without a ground mark. Choose
a vertex $s_0$ in such a component and assign it weight
$\lambda_{s_0}=1\in\mathbb Q$. For an edge with endpoint
coefficients $a,b\ne0$, extend the weights along the unique
tree path by requiring
\[
 a\lambda_s+b\lambda_t=0.
\]
Every vertex is connected to $s_0$ by a finite path, so this
assigns a nonzero rational weight to every vertex, with no
consistency ambiguity. Assign weight zero on all other
components. The resulting nonzero linear functional
\[
 \ell:\mathbb Q^{(S)}\longrightarrow\mathbb Q,
 \qquad \ell(s)=\lambda_s,
\]
vanishes on the image of $\delta\otimes\mathbb Q$.
This contradicts its surjectivity.

It follows that every component is either a tree with
exactly one ground mark, or a component with exactly one
finite cycle and no ground marks. In the second case,
removing the cycle edges leaves trees attached to its
vertices; an additional connection would create another
cycle. The finite exponent matrix of the cycle is square
and nonsingular. Indeed, its columns are the cycle
relators, all supported on the same number of cycle
vertices, and injectivity of $\delta$ makes these columns
linearly independent over $\mathbb Q$.

\paragraph{Solving the entire presentation.}
For a component with a ground mark, solve that one-variable
equation in $E_H$, with its coefficients read through $j$.
Its exponent is nonzero, so this is a finite nonsingular
system. For a component with a cycle, solve simultaneously
its finitely many cycle equations in $E_H$, using the
square nonsingular matrix just proved. Make these choices
for all components.

We now proceed by finite distance from the chosen root
vertex or cycle. Once the value at a parent vertex is
assigned, its edge to a child gives a single equation in
that child variable. Every other letter in the edge word
is either a coefficient from $H$ or this already assigned
parent variable. Substitution preserves the nonzero
exponent sum of the child. The equation therefore has a
solution in $E_H$. Assign such values simultaneously at
each distance. There is only one parent edge for each
new vertex, and every vertex has finite distance from
its root or cycle. Thus all variables are assigned and
all relators are satisfied. Arbitrary valence and
arbitrary cardinalities of the set of components cause
no change in this recursion.

The homomorphism $H*F(S)\to E_H$ defined by $j$ and these
variable values kills every specified $w_r$. Its descent
is the promised $\Phi:G\to E_H$. The equality of its
restriction with $j$ proves the coefficient injection
as well.

\paragraph{The countable-core consequence.}
Fix $g\ne1$ in the group of the final assertion. Apply
the countable homology-core construction in A.19 to a
countable subcomplex containing the loop representing
$g$, all cells in $S_*$ and $R_*$, and their attaching
supports. It produces an actual countable subcomplex
$Y\subset X$ whose inclusion is injective on fundamental
groups and an isomorphism on all ordinary integral
homology. Since $X$ has one vertex, so does $Y$.
The actual relative boundary
\[
 C_2(X,Y;\mathbb Z)\longrightarrow C_1(X,Y;\mathbb Z)
\]
is an isomorphism, as in (A.19.4). Write
$H=\pi_1(Y)$, which is countable.

The relative attaching words for the remaining two-cells
are the original words with the letters of $Y$ read as
coefficients in $H$. The exceptional relators lie in
$Y$. Every other word has at most two remaining
variables, since $S_*\subset Y$, and each remaining
variable retains its nonzero exponent sum. This is
exactly a presentation of the form (A.25.1), so its
homomorphism $G\to E_H$ is injective on $H$. It sends
$g$ to a nonidentity element. Restricting its codomain
to its countable image gives the required countable
quotient.
\end{proof}

The consequence applies, in particular, to an anchored Camm-type
word pattern
\[
 a^{n_r}p_{\alpha(r)}^{u_r}
       =b^{m_r}q_{\beta(r)}^{v_r},\qquad u_rv_r\ne0,
\]
with the anchors $a,b$ in the countable exceptional set of
generators, whenever the specified two-dimensional model has
the small homology in A.25. The proof constructs its countable
detecting homomorphisms. A modification with more variables
or with variable occurrences of exponent sum zero requires
a further argument; no conclusion for that modification
is claimed here.
